\section{Localization of the original pointwise forcing}
\label{sec:v96-forcing}
Take $H=\Pi$ in the occupation tests. Partition its nonadmitted targets
into the following disjoint Borel classes. The first, $N_T$, has no
available comparison interval of length $r$ in either direction. The
second, $N_D$, has at least one available interval but no available sign
with $d_{e,r}\le D$. The third, $N_O$, has such a bounded-compression
sign but no such sign with $o_{\Pi,e,r}\ge\beta$. Thus physical
interruption, compression and insufficient same-label occupation are
charged separately and only once. Define their original weighted
source sum by
\begin{equation}\label{eq:v96-failure-source}
 Z_{r,D,\beta}^\varepsilon(t)
  =\int_{F=t}W^\varepsilon
           (\one_{N_T}+\one_{N_D}+\one_{N_O})\,d\mu_t.
\end{equation}
In particular $0\le Z_{r,D,\beta}^\varepsilon\le b^\varepsilon$.
No new source is placed on a comparison-domain boundary.

\begin{theorem}[Prescribed-radius localization of the full source]
\label{thm:v96-pointwise}
At every finite count and almost every roof,
\begin{equation}\label{eq:v96-pointwise}
 b^\varepsilon(t)\le Z_{r,D,\beta}^\varepsilon(t)
 +\frac{D}{\beta r}\int_{t-r}^{t+r}
             \{b^\varepsilon(u)+\kappa_\varepsilon(r)p(u)\}\,du.
\end{equation}
The same bound dominates the bidirectional forcing of module 228,
because that forcing is bounded by the original target source.
Consequently, with $\varepsilon=\varepsilon(B)$,
\begin{equation}\label{eq:v96-endpoint-bound}
 \mathcal E_M\le C_MB^{-1/192}
   +\mathcal H_M(Z_{r,D,\beta}^\varepsilon)
   +C\frac{D}{\beta}
                \{\varepsilon^{1/16}+\varepsilon^{-1}\sqrt r\}.
\end{equation}
The first and last terms are estimated. The middle term retains the
complete weighted mass of the three failed geometric tests.
\end{theorem}
\begin{proof}
The exact partition of $b^\varepsilon$ into admitted and nonadmitted
parts, followed by Theorem~\ref{thm:v96-anchor}, proves
\eqref{eq:v96-pointwise}. It is a statement at the target roof, not
an integral in the target variable. The retained bidirectional
forcing obeys $0\le J\le b^\varepsilon$ at finite count, so no
minimum is moved inside the physical-sheet integral.

For fixed $B,r,D,\beta$ the complete source average is at most
$C\varepsilon^{1/16}$ in the normalized limit. The unweighted
average has a common bound by the same theorem at auxiliary guard
parameter one, where $W^1=1$. Apply these two estimates to
\eqref{eq:v96-pointwise}, then insert the original positive scalar
remainder. This proves \eqref{eq:v96-endpoint-bound}, without any
spectral estimate outside the exact event.
\end{proof}

\begin{corollary}[A legal outer schedule and its exact obligation]
\label{cor:v96-schedule}
Suppose $D(B)/\beta(B)\le C\varepsilon(B)^{-\tau}$, where
$0\le\tau<1/16$, and set $r(B)=c\varepsilon(B)^N$ with
$N>2(1+\tau)$ and sufficiently small fixed $c>0$. The original
pointwise theorem is equivalent to
\begin{equation}\label{eq:v96-failure-zero}
 \mathcal H_M(Z_{r(B),D(B),\beta(B)}^{\varepsilon(B)})
                           \longrightarrow0\quad(B\to\infty)
\end{equation}
for every fixed $M$. The same schedule preserves the equivalence
with vanishing of the retained bidirectional forcing.
\end{corollary}
\begin{proof}
The explicit errors have powers $\varepsilon^{1/16-\tau}$ and
$\varepsilon^{N/2-1-\tau}$, both positive. This proves sufficiency
from \eqref{eq:v96-endpoint-bound}. Conversely the original
pointwise law and the positive scalar remainder give
$\mathcal H_M(b^{\varepsilon(B)})\le C_MB^{-1/192}$.
The domination $Z\le b^\varepsilon$ proves necessity. The retained
bidirectional equivalence applies to this same prescribed radius,
since $(1+\varepsilon^{-1})r\to0$. One simple schedule is
$D=2$, $\beta=1/4$, $r=c\varepsilon^4$. All these parameters,
including each sampling radius, are fixed before their collision
limit; auxiliary bands are enlarged afterwards and $B$ is last.
\end{proof}

For an original exceptional selector $H_{a,\delta}$ the identical
three-way partition is made using returns to $H_{a,\delta}$ itself.
Theorem~\ref{thm:v96-source-height} then localizes its essential-height
charge to failed continuation, compression or selected-event occupation,
with the explicit admitted error $C(D/\beta)(a^\alpha+\delta^\gamma)$.
Together with Theorem~\ref{thm:v96-packet-decay}, this controls the
recurrent part of the two-nongrazing packet and the complete recurrent
radial core, not just a finite annulus head.

No theorem here proves \eqref{eq:v96-failure-zero} for the full
Lorentz source. A short two-nongrazing pulse may have too little
same-label occupation, and a critical target may fail physical
continuation or the compression test. These possibilities remain
positive original-source terms. In the elementary comparison model
$\widehat A\equiv1$ and $\Pi=(0,\delta)$, occupation over a fixed
one-sided radius is at most $\delta/r$; the theorem does not erase
that pulse. This records the boundary between a proved persistence
mechanism and the still required exhaustion of all spike mechanisms.
The original scalar, arithmetic and unrestricted downstream claims
are therefore not promoted.
